 \section{Explantation of a Spanned Slice}

  \ctt{Here we follow the original proof, but change the triminalogy, $m$ will be the number of leaves in the explanation, then in the end we conclude that the number of faults is at least $m/\Delta$. So $|R| \le f(\frac{m}{\Delta})$. }

  \begin{definition}
    Consider a Tooms contour graph induced by Toric encoded circuit, with sweep-cycle parameter $\xi = 8$, and  let $\alpha, \beta$ be constants and $\mathcal{X}_{0}, \mathcal{X}_{1}, \mathcal{X}_{2}, \mathcal{X}_{3} \in \mathbb{R}$ be quadruple, defining such:
    \begin{equation*}
      \begin{split}
      \alpha &\colonequals 10(d+1)^{2}, \ \ \  \beta \colonequals 5 (d+1) \\
      \mathcal{X}_{r} &\colonequals 1 - r \frac{d+1}{\alpha} = 1 -  \frac{r}{10(d+1)}
      \end{split}
    \end{equation*}
    Given sets $W$ of points, $R$ of arrows and $F$ of forks, we define $U$ as the set of vertices of the graph induced by the edges of $R \cup F$. The sets $W,R$ and $F$ form an explanation of a spanned slice $\hat{K} = \expval*{ K, v_{1}, .., v_{d+1}}$ of order $\text{order}(\hat{K}) = r$, if and only if: 
    \begin{enumerate}
      \item $W \cup \{v_{1} ,.., v_{d+1}\} \subset U \subset \History{K}$ and the graph $(U, R\cup F)$ is connected. 
      \item All elements of $W$ are leaves. 
      \item For $m = |W|$ and $s=\Size{K}$ we have $|R| \le \alpha (m - \mathcal{X}_{r}) - \beta s $ and $|F| < m -1 $. 
    \end{enumerate}
  \end{definition}

  \begin{claim}
    \label{claim:edgesin}
    Every spanned slice $K = \expval*{K,v_{1},..,v_{d+1}}$ has an explanation.
  \end{claim}

  \begin{proof}
    Let $h = \Time{K} - \min_{u \in \History{K}}\Time{u}$. The proof is by induction on $h$. 
    \begin{enumerate}
      \item Base: $h=0$, meaning $T$ is constant on $\History{K}$, so no arrows connect points of $\History{K}$. Thus, $\History{K}$, as well as $\mathbf{K}$, consists only of a single leaf, namely $\hat{K}$ is a result of a fault, and $r = \text{order}(\mathcal{X}_{r}) \in \{0,2\}$. In that case, there is a single vertex $u \in V_{\mathcal{H}}$ such that $U = W = \{u\}$, so $m=1$, $s=0$, and $|R|, |F| = 0$. It is easy to check that those values satisfy the requirement. 
      \item Induction Step: $h \neq 0$ implies $K \neq \History{K}$. Thus, by \cref{claim:prevslices}, we know that for some $k$ there exist spanned slices $\hat{K}_{0}, .., \hat{K}_{k}$ and forks $F_{1},..,F_{k}$ that satisfy \cref{claim:prevslices}. In particular, for $r < 3$ it holds that: 
        \begin{equation*}
          \begin{split}
            \sum_{i=0}^{k}\Size{\hat{K}_{i}} +  \sum^{k}_{i=1}\Size{F_{i}} \ge \Size{\hat{K}}  
          \end{split}
        \end{equation*}
        Whereas, for $r=3$: 
        \begin{equation*}
          \begin{split}
            \sum_{i=0}^{k}\Size{\hat{K}_{i}} +  \sum^{k}_{i=1}\Size{F_{i}} \ge \Size{\hat{K}} + \xi = \Size{\hat{K}} + 8
          \end{split}
        \end{equation*}
        Consider $i \in \{0, \ldots, k\}$, and denote $s_{i} = \Size{\hat{K}_{i}}$. By the inductive hypothesis and since $\Time{K_{i}} = \Time{K} - 1$, we know that $\hat{K}_{i}$ has an explanation formed by a set of leaves $\mathbf{W}_{i}$, a set of arrows $\mathbf{R}_{i}$, and a set of forks $\mathbf{F}_{i}$. Let $m_{i} = |\mathbf{W}_{i}|$. Thus, the sets $\mathbf{R}_{i}$ and $\mathbf{F}_{i}$ have at most $\alpha (m_{i} - \mathcal{X}_{r-1}) - \beta s_{i}$ arrows and $m_{i} - 1$ forks. We define:
        \begin{enumerate}
          \item $\mathbf{W} = \cup_{i=0}^{k} \mathbf{W}_{i}$, Notice that $\mathbf{W}_{i} \subset K_{i}$ and that $\Time{K_{i}} = \Time{K_{j}}$ for all $i,j$, thus by \Cref{claim:disj} $\mathbf{W}$ is a union of disjoint set, Hence $W$ contains  $m = \sum_{i=0}^{k}m_{i}$ leaves. 
          \item $\mathbf{F} = \{ F_{1},..,F_{k} \} \cup_{i = 0}^{k}\mathbf{F}_{i}$ with at most $ k + \sum_{i=0}^{k}{m_{i} -1 } = m- 1$ elements. 
          \item $\mathbf{R} = \{ \{v_{i}, \prei{v_{i}} \} \} \cup_{i=0}^{k}{ \mathbf{R}_{i} }$. For bounding the number of arrows in $\mathbf{R}$ we condition upon whether $r = 3$, For convenient, let us denote by $Y$ the indicator, which indicate if $r = 3$.   
            \begin{equation*}
              \begin{split}
                & d+1 + \sum^{k}_{i=0}{\alpha (m_{i}-\mathcal{X}_{r-1})  - \beta s_{i} } \\ 
                \le & d+1 + \alpha m - \mathcal{X}_{r-1}\alpha(k+1) - \beta\left(s + Y\xi - k \cdot \Size{\text{fork}}  \right) \\
                = & \alpha \left(m - \mathcal{X}_{r}\right) - \beta s + \left(\alpha ( \mathcal{X}_{r} - \mathcal{X}_{r-1} ) + d+1 -  \mathcal{X}_{r-1} k \alpha - \beta Y \xi + \beta k \cdot \Size{\text{fork}} \right)
              \end{split}
            \end{equation*}
So it's left to show that the term in the closure is negative for all integers $k$:
            \begin{equation*}
              \begin{split}
\alpha ( \mathcal{X}_{r} - \mathcal{X}_{r-1} ) + d+1 -  \mathcal{X}_{r-1} k \alpha - \beta Y \xi + \beta k \cdot \Size{\text{fork}} 
              \end{split}
            \end{equation*}
            We do that by showing that for the choice of the parameter values, $\alpha, \beta, \mathcal{X}_{0}, \ldots, \mathcal{X}_{3}$, the slope is negative, and then the value of the term at $k$ equals zero is negative.  For the slope to be negative, we require: 
            \begin{equation*}
              \begin{split}
                - \alpha \mathcal{X}_{r-1}  +  \beta \Size{\text{fork}} \le 0 & \Rightarrow   \Size{\text{fork}} \le \frac{\alpha}{\beta} \mathcal{X}_{r-1} 
              \end{split}
            \end{equation*}
            For negativity at the origin, $k=0$, we require for $r < 3$:
            \begin{equation*}
              \begin{split}
                (\mathcal{X}_{r}-\mathcal{X}_{r-1})  & \alpha  + d + 1  \le 0 
                \\ &  \Rightarrow \mathcal{X}_{r} \le  \mathcal{X}_{r-1} - \frac{d+1}{\alpha}  
  \end{split}
\end{equation*}
And for $r =3$:
\begin{equation*}
  \begin{split}
(\mathcal{X}_{3}-\mathcal{X}_{2}) & \alpha  + d + 1  -\beta \xi \le 0  \\ 
& \Rightarrow \mathcal{X}_{3} \le  \mathcal{X}_{0} - 3\frac{d+1}{\alpha} + \frac{\beta}{\alpha}\xi  
              \end{split}
            \end{equation*}
So in total: 
\begin{equation*}
  \begin{split}
&  \mathcal{X}_{1} \le  \mathcal{X}_{0} - \frac{d+1}{\alpha}  \\
&  \mathcal{X}_{2} \le  \mathcal{X}_{0} - 2\frac{d+1}{\alpha}  \\ 
&  \mathcal{X}_{3} \le  \mathcal{X}_{0} - 3\frac{d+1}{\alpha} + \frac{\beta}{\alpha}\xi  \\
&  \mathcal{X}_{0} \le  \mathcal{X}_{0} - 4\frac{d+1}{\alpha}  + \frac{\beta}{\alpha}\xi  \\
 & \Size{\text{fork}} \le \frac{\alpha}{\beta} \mathcal{X}_{r} 
  \end{split}
\end{equation*}
The first three inequalities hold by the definitions of the $\mathcal{X}_{r}$, the fourth is satisfied since $\frac{\beta}{\alpha}\xi = 4 \frac{d+1}{\alpha}$, the last is obtained by the fact that the maximal size of a fork equals $d$ (which matches the fork between two extremums of a face), and that $\frac{\alpha}{\beta}\mathcal{X}_{r}$ is greater than $2d$.
        \end{enumerate}
         Finally, by \cref{claim:prevslices}, the graphs induced by $\mathbf{R}_{i} \cup \mathbf{F}_{i}$ are connected, thus the graph induced by $\mathbf{R}\cup \mathbf{F}$ is also connected. Therefore, $\mathbf{W}, \mathbf{R}, \mathbf{F}$ form an explanation of $K$. 
    \end{enumerate}
    \ctt{ Check what happened if we set $\xi = \sqrt{d}, \beta = \sqrt{d}$ is that better? }
  \end{proof}



